\documentclass[10pt,a4paper]{amsart}
\usepackage[margin=19mm]{geometry}
\usepackage{amsmath,amssymb,mathtools}
\usepackage[scr=rsfs,cal=euler]{mathalfa}
\usepackage{xfrac}
\usepackage[unicode,hidelinks,bookmarks=false]{hyperref}
\newcommand{\ie}{i.e.\ }
\newcommand{\cf}{cf.\ }
\newcommand{\resp}{respectively\ }
\newcommand{\Subsection}{\subsection}
\providecommand{\nobreakdash}{\nobreak\hbox{-}}
\setlength{\emergencystretch}{2em}
\multlinegap0pt
\usepackage{bm}
\usepackage{bbm}
\usepackage{dsfont}
\usepackage{xspace}
\usepackage{cancel}
\usepackage{mathtools}
\usepackage{amsthm}
\makeatletter
\@ifundefined{theorem}{\newtheorem{theorem}{Theorem}}{}
\@ifundefined{lemma}{\newtheorem{lemma}[theorem]{Lemma}}{}
\@ifundefined{proposition}{\newtheorem{proposition}[theorem]{Proposition}}{}
\makeatother
\newcommand{\cH}{\mathcal{H}}
\title[A 60-Vertex Lower Bound for Cubic Bipartite Counterexamples ]{A 60-Vertex Lower Bound for Cubic Bipartite Counterexamples to the Erd\H{o}s-Gy\'arf\'as Conjecture}
\author{Julius Tranquilli}
\date{Source: arXiv:2608.02675v1 (CC BY 4.0)}
\begin{document}
\maketitle

\noindent\emph{Source and licence.} A 60-Vertex Lower Bound for Cubic Bipartite Counterexamples to the Erd\H{o}s-Gy\'arf\'as Conjecture. Version arXiv:2608.02675v1 (CC BY 4.0), distributed under the Creative Commons Attribution 4.0 International licence, as recorded in the arXiv licence metadata for this exact version and shown on the abstract page https://arxiv.org/abs/2608.02675v1. Full rights notice: licenses/arxiv-2608.02675-CC-BY-4.0.txt.\par\smallskip

\noindent\textit{Editorial scope.} Counted unit: the Proposition whose source label is \texttt{prop:triangle-coverage} in the article (source0.tex lines 458--463, source label \texttt{prop:triangle-coverage}), delivered together with the article's own complete original proof of it (source0.tex lines 465--496, one \texttt{proof} environment of 32 lines). No mathematics is written, repaired or re-derived: statement and proof are the article's own text line for line. Required context retained uncounted: lem:triangle-orbits (lines 369--380). Every definition, display and label of those lines is retained unchanged. Omitted from this package and not redistributed: 1--368, 381--457, 464--464, 497--866. Nothing inside a retained excerpt is omitted or altered: every retained body is delivered exactly as the article's lines are written, so no omission receipt of any kind accompanies this package. Citations: the retained excerpts carry no citation command at all, so no bibliography entry of the article is transcribed and no reference is redistributed. Reference adaptation: none was needed and none was made; every reference inside the delivered unit resolves inside it, so no unresolved reference or citation is produced. Printed slips kept literal: no printed word, symbol, hypothesis or number of the counted statement or of its proof was altered. Layout and class: the article's own class is \texttt{article} with packages including fontenc, inputenc, geometry, amsmath, amssymb, amsthm, booktabs, tabularx, xcolor, tikz, microtype, hyperref; the wrapper is the lane's pinned \texttt{amsart} editorial wrapper at 10pt on A4, which restates the theorem-like environments the retained text needs and adds the article's own macro definitions that the retained text needs (\texttt{\textbackslash cH}), with extra wrapper packages (bm, bbm, dsfont, xspace, cancel, mathtools). The delivered numbering is the unit's own. Assets: no retained span contains a figure environment, a graphics-inclusion command, a PDF-page inclusion, a table or a TikZ picture; no image file is copied into this package, the package declares no asset dependency and no image file is redistributed with it. Licence: version arXiv:2608.02675v1 of this article is distributed under the Creative Commons Attribution 4.0 International licence, as recorded in the arXiv licence metadata for this exact version and shown on the abstract page https://arxiv.org/abs/2608.02675v1. A full rights notice is delivered at licenses/arxiv-2608.02675-CC-BY-4.0.txt. Characters: every retained excerpt is pure ASCII, so the delivered engine, XeLaTeX with its default Latin Modern text font, reports no missing character.
\begin{lemma}[Triangle-root orbits]\label{lem:triangle-orbits}
Let \(\cH\) be a linear \(3\)-uniform configuration containing a Berge
triangle. After relabeling, its three triangle blocks are
\[
  \{0,1,3\},\qquad \{1,2,4\},\qquad \{0,2,5\}.
\]
Up to the stabilizer of this rooted triangle, the final block through point
\(0\) is one of
\[
  \{0,4,6\},\qquad \{0,6,7\}.
\]
\end{lemma}

\begin{proposition}[Triangle-rooted coverage]\label{prop:triangle-coverage}
Every connected linear symmetric \(v_3\)-configuration with \(v\leq29\)
that contains a Berge triangle and has no Berge cycles of lengths \(4\) or
\(8\) occurs as a completed state in one of the two triangle-rooted
cap-\(29\) restricted-growth trees.
\end{proposition}

\begin{proof}
Choose a Berge triangle and apply Lemma~\ref{lem:triangle-orbits}. Install
its three blocks and the appropriate fourth block through point \(0\).
The preinstalled block \(\{1,2,4\}\) is lexicographically first among the
blocks with least point \(1\): any other such block cannot reuse a point
already paired with \(1\). Begin the ordinary restricted-growth recursion at
point \(1\), after this block.

Maintain the following invariant before processing a point \(p\): every
target block with smaller least point has been inserted; the inserted target
blocks through \(p\) form a lexicographic initial segment; and introduced
labels are consecutive and follow first occurrence. The normalized root
establishes the invariant at \(p=1\).

Let \(B\) be the first target block through \(p\) not yet inserted. Every
block containing \(p\) and a smaller point was inserted when that smaller
point was processed, so the other two points of \(B\) have labels greater
than \(p\). Previously unseen points receive the next one or two labels.
Thus \(B\) occurs among the generator's proposals and is later than the
preceding block through \(p\). It passes the degree and pair tests because
the target is linear, and it passes the cycle tests because the target has no
Berge cycles of lengths \(4\) or \(8\). Inserting \(B\) preserves the
invariant; when \(p\) becomes cubic, the recursion advances to the least
unfinished introduced point.

If the introduced points formed a proper subset closed under their incident
blocks, the incidence graph would be disconnected. Hence connectedness
ensures that every target point is eventually introduced. Since there are at
most \(29\) points, no label outside the cap is needed. Once all introduced
points are cubic, the tree recognizes the completed configuration
immediately, even if fewer than \(29\) points were introduced.
\end{proof}

\end{document}
